\documentclass[12pt]{article}

\usepackage[paper=letterpaper,margin=2.5cm]{geometry} % Set Margins

%% Math and math fonts
\usepackage{amsmath, amsthm, amssymb, amsfonts}
\usepackage{bbold} % for \mathbbm{1}

% date
\usepackage[nodayofweek]{datetime}

% Color
\usepackage{color, xcolor}

% Misc
\usepackage{environ}  % \collect@body in asmmath
\usepackage{graphicx} % \includegraphics options
\usepackage{mdframed} % text boxes
\usepackage{indentfirst} % Indent first paragraph after section header
\usepackage[shortlabels]{enumitem} % Control enumerate items with [(a)]
\usepackage{comment} % Comments
\usepackage{fancyhdr} % Headers and footers

% Tables
\usepackage{array}

% Sub-figures and figure placement
\usepackage{caption}
\usepackage{subcaption}
\usepackage{float} 

% Graphing
\usepackage{pgfplots}
\pgfplotsset{compat=1.17}
\usepackage{tikz}

% Title Placement
\usepackage{titling}
\setlength{\droptitle}{-6em}
\setlength{\headheight}{15pt}

%set indent to 
\setlength{\parindent}{0pt}
\setlength{\parskip}{10pt}

% Hyper refs
\usepackage{hyperref}
\hypersetup{
    colorlinks=true,
    linkcolor=blue,
    urlcolor  = blue,
    filecolor=magenta,      
    urlcolor=blue,
    citecolor = blue,
    anchorcolor = blue
}

% % Citation management
\usepackage{natbib}
\bibliographystyle{abbrvnat}
\setcitestyle{authordate,open={(},close={)}}

% % Move the header up
\usepackage{titling}
%\posttitle{\par\centering}
\preauthor{\centering\normalsize}
\postauthor{\par\centering\vspace{-1em}}
\predate{\centering\normalsize}
\postdate{\par\centering\vspace{-1em}}

% ----------------------------------------
% TITLE
% ----------------------------------------

\pagestyle{fancy}

\lhead{Camacho}
\chead{Class 8}
\rhead{FAME}

\title{FAME Class 8 Notes -- The Math of Environmental Change \& Aggregation and Comparisons in Environmental Systems}
\author{Mariana Camacho}
\date{November 19, 2026}

\begin{document}
\maketitle
\section{Revisiting the implicit function theorem}
Imagine that we have the following function: \(m=x^2+y^2\)

\begin{figure}[h!]
\centering

\begin{tikzpicture}[scale=1.15]

    % Y-axis: Latitude
    \draw[thick] (0,0) -- (0,6.5)
        node[above] {Latitude};

    % X-axis: Longitude
    \draw[thick] (0,0) -- (8,0)
        node[right] {Longitude};

    % 100 m curve
    \draw[thick]
        plot[smooth, tension=0.8] coordinates {
            (0.4,5.2)
            (0.7,3.9)
            (1.5,2.5)
            (2.6,1.5)
            (4.0,0.9)
            (5.0,0.7)
        };
    \node[right] at (5.05,0.7) {$100\,m$};

    % 120 m curve
    \draw[thick]
        plot[smooth, tension=0.8] coordinates {
            (1.3,5.6)
            (1.6,4.5)
            (2.4,3.2)
            (3.5,2.3)
            (4.7,1.8)
            (5.4,1.7)
        };
    \node[right] at (5.45,1.7) {$120\,m$};

    % 140 m curve
    \draw[thick]
        plot[smooth, tension=0.8] coordinates {
            (2.7,6.0)
            (2.9,5.0)
            (3.6,3.9)
            (4.6,3.1)
            (5.7,2.7)
            (6.3,2.7)
        };
    \node[right] at (6.35,2.7) {$140\,m$};

    % 160 m curve
    \draw[thick]
        plot[smooth, tension=0.8] coordinates {
            (4.0,6.3)
            (4.2,5.6)
            (4.8,4.9)
            (5.6,4.5)
            (6.5,4.4)
            (7.0,4.5)
        };
    \node[right] at (7.05,4.5) {$160\,m$};

\end{tikzpicture}

\caption{Contour lines by elevation (m).}
\label{fig:contour_lines}

\end{figure}

We can think of these as contour lines for elevation. Now, let's take the \textbf{total derivative}. 

\(\mathrm{d}M =
\frac{\partial M}{\partial x}\mathrm{d}x
+
\frac{\partial M}{\partial y}\mathrm{d}y = 2x \partial x + 2y \partial y =
\begin{bmatrix}
2x & 2y
\end{bmatrix}
\begin{bmatrix}
 \partial x\\
\partial y
\end{bmatrix}\)

Let's focus on the \textbf{partial derivative} with respect to x:

\(\frac{\partial M}{\partial x}=2x\)

We can rewrite this as:

\(\partial M=2x \partial x\)

Now, let's look at another example:

\(F= \frac{x}{y}=xy^{-1}\)

Now, let's differentiate and set it equal to zero.

\(\mathrm{d}F=y^{-1} \partial x + (-xy^{-2}) \partial y = 0\)

We can rewrite this as:

\(y^{-1} \partial x +=xy^{-2} \partial y\)

and consequently:

\( \frac{y^2}{xy}= \frac{y}{x} = \frac{y^{-1}}{xy^2} = \frac{\partial y}{\partial x}\)

\section{An example from ecology}
Now, let's think about a different problem. Ecologists have found that the number of species found in a habitat depends on the area. 

\(S(A)=20A^{0.25}\)

The reserve is currently \(10,000 \ km^2\), but policymakers are analyzing several alternatives. 

We want to know the \textbf{sensitivity} of the species count to reserve area. When you hear the word sensivity, you should immediately think of a derivative!

In this case, the derivative will be:

\( \frac{dS(A)}{dA}=0.25*20A^{-0.75}\)

This relationship directly follows from \(S(A)\), so if we don't think that this is right, we might want to go back and reexamine \(S(A)\). In practice, sometimes it's easier to test out \( \frac{dS(A)}{dA}\) than \(S(A)\). 

We can also write the derivative as:

\( \frac{S(A+t)-S(A)}{\epsilon} \)

We might also be interested in the percentage change. 

\( \frac{\frac{dS(A)}{dA}}{S}= \% \Delta\)

We can write this as:

\( \frac{\partial ln(S(A)}{\partial A}= \frac{1}{S(A)} \frac{\partial S}{ \partial A} \)

We can call these \textbf{elasticities}. 

\section{Introduction to aggregation}

Integration let's us think about \textbf{areas}. How could we approximate the area between a and F of the function below?

\begin{figure}[h!]
\centering

\begin{tikzpicture}[scale=1.15]

    % Axes
    \draw[thick] (0,0) -- (0,6)
        node[above] {$y$};

    \draw[thick] (0,0) -- (8,0)
        node[right] {$x$};

    % Function curve
    \draw[thick]
        plot[smooth, tension=0.8] coordinates {
            (0,2.5)
            (0.8,3.4)
            (1.7,4.2)
            (2.5,4.8)
            (3.5,5.25)
            (4.7,5.55)
            (5.8,5.7)
            (6.5,5.75)
        };

    % Vertical line at a
    \draw[thick] (2.5,0) -- (2.5,4.8);

    % Vertical line at F
    \draw[thick] (5.8,0) -- (5.8,5.7);

    % Horizontal dashed line indicating f(a)
    \draw[dashed, thick] (0,4.8) -- (2.5,4.8);

    % Labels
    \node[left] at (-0.05,4.8) {$f(a)$};
    \node[below] at (2.5,0) {$a$};
    \node[below] at (5.8,0) {$F$};

\end{tikzpicture}

\caption{Illustration of an increasing function.}
\label{fig:increasing_function}

\end{figure}

Our initial approximation might be:

\(F(T)*(T-a)\)

We could also try:

\(F(a)*(T-a)\)

We can do better if we use a \textbf{Riemann sum} - make small rectangles and add them all together, using the \(\sum \) operator.

We can also introduce continuous notation: \( \int \, dx \)

We can write the area as:

\(A(x)=f(x) \Delta x \rightarrow f(x) \epsilon\)

We can also think of:

\( \frac{A(x+\varepsilon)-A(x)}{\epsilon} = \frac{ \partial A}{ \partial x} \rightarrow \int_a^T \mathrm{d}A = \int_a^T A'(x)\,\mathrm{d}x \)

We can rewrite this as:

\(A = \int_a^T A'(x)\,\mathrm{d}x\)

\section{An example from Arizona}

This is the current layout of a property in Arizona. 

\begin{figure}[h!]
\centering

\begin{tikzpicture}[scale=1.15]

    % Outer property boundary
    \draw[thick]
        (0,0.5) --
        (0.3,5.8) --
        (2.0,6.1);

    \draw[thick]
        (2.0,6.1) --
        (7.1,6.4) --
        (8.5,6.3) --
        (8.5,0.2) --
        (6.5,0.4);

    \draw[thick]
        (6.5,0.4) --
        (6.5,4.2);

    \draw[thick]
        (0,0.5) --
        (4.7,0.1);

    % House
    \draw[thick]
        (2.0,6.1) --
        (2.1,4.1) --
        (6.5,4.4) --
        (6.5,6.4);

    % House label
    \node at (4.3,5.3) {\large House};

    % Driveway
    \draw[thick]
        (4.7,0.1) --
        (5.4,4.3);

    % Driveway label
    \node[rotate=70] at (5.7,1.8) {Drive};

    % Grass labels
    \node at (2.4,2.5) {Grass};
    \node at (7.5,2.8) {Grass};

\end{tikzpicture}

\caption{Property layout.}
\label{fig:property_layout}

\end{figure}

We want to reduce water consumption on grass watering, so we have the following plan:


\begin{figure}[h!]
\centering

\begin{tikzpicture}[scale=1.15]

    % =================================
    % OUTER PROPERTY BOUNDARY
    % =================================

    \draw[thick]
        (0,0.5) -- (0.3,5.8) -- (2.0,6.1);

    \draw[thick]
        (2.0,6.1) -- (7.1,6.4)
        -- (8.5,6.3) -- (8.5,0.2)
        -- (6.5,0.4);

    \draw[thick] (6.5,0.4) -- (6.5,4.4);

    \draw[thick] (0,0.5) -- (4.7,0.1);

    % =================================
    % HOUSE
    % =================================

    \fill[gray!15]
        (2.0,6.1) -- (2.1,4.1)
        -- (6.5,4.4) -- (6.5,6.4) -- cycle;

    \draw[thick]
        (2.0,6.1) -- (2.1,4.1)
        -- (6.5,4.4) -- (6.5,6.4);

    \node at (4.3,5.3) {\large House};

    % =================================
    % GRASS AREA
    % =================================

    \fill[green!15]
        (2.1,4.1)
        -- (5.4,4.3)
        -- (4.7,0.1)
        -- (2.5,0.3)
        .. controls (1.6,0.6) and (1.3,1.8) .. (1.5,2.6)
        .. controls (1.6,3.3) and (1.9,3.8) .. (2.1,4.1)
        -- cycle;

    % Curved left boundary of grass
    \draw[thick]
        (2.1,4.1)
        .. controls (1.9,3.8) and (1.6,3.3) .. (1.5,2.6)
        .. controls (1.3,1.8) and (1.6,0.6) .. (2.5,0.3);

    % Sloping right boundary of grass
    \draw[thick] (5.4,4.3) -- (4.7,0.1);

    % =================================
    % DRIVEWAY
    % =================================

    \fill[gray!25]
        (5.4,4.3) -- (6.5,4.4)
        -- (6.5,0.4) -- (4.7,0.1) -- cycle;

    \node[rotate=80] at (5.9,2.1) {Drive};

    % =================================
    % SHRUBS: RIGHT-HAND STRIP
    % =================================

    \fill[purple!15]
        (7.1,6.4) -- (8.5,6.3)
        -- (8.5,0.2) -- (6.5,0.4)
        -- (6.5,4.4) -- (7.1,6.4) -- cycle;
    
    \draw[thick]
        (7.1,6.4) -- (8.5,6.3)
        -- (8.5,0.2) -- (6.5,0.4);

    % =================================
    % TREES: LEFT-HAND STRIP
    % =================================

    \foreach \x/\y in {
        0.7/1.0,
        0.8/2.0,
        0.7/3.0,
        0.8/4.0,
        0.9/5.0
    } {
        \draw[green!40!black, thick, fill=green!35]
            (\x,\y) circle (0.22);
    }

    % =================================
    % LABELS
    % =================================

    \node at (3.5,2.4)
        {\textcolor{green!40!black}{Grass}};

    \node[rotate=90] at (7.9,3.3)
        {\textcolor{purple!70!black}{Shrubs}};

    \node at (1.0,5.6)
        {\textcolor{green!40!black}{Trees}};

\end{tikzpicture}

\caption{Proposed landscape layout.}
\label{fig:landscape_layout}

\end{figure}

How should we go about measuring the areas so we can do this with the lowest cost possible? 

The shrubs are easy, we can use a rectangle. The trickiest part is the curved part on the bottom left, in the boundary between the trees and the grass. We might think of approximating it as:


\begin{figure}[h!]
\centering

\begin{tikzpicture}[scale=1.15]

    % Square boundaries
    \draw[thick] (0,0) -- (0,5);
    \draw[thick] (0,5) -- (5,5);
    \draw[thick] (5,5) -- (5,0);
    \draw[thick] (0,0) -- (5,0);

    % Increasing, convex function
    \draw[thick]
        plot[smooth, tension=0.8] coordinates {
            (0,0)
            (1.0,0.08)
            (2.0,0.35)
            (3.0,1.0)
            (4.0,2.2)
            (4.6,3.5)
            (5.0,5.0)
        };

    % Origin
    \node[below left] at (0,0) {$0$};

    % Axis labels
    \node[above left] at (0,5) {$f$};
    \node[below right] at (5,0) {$T_x$};

\end{tikzpicture}

\end{figure}

We can define:

\(y=ax^2+bx\)

Taking the integral:

\(A(x) = \int_0^T (ax^2 + bx)\,\mathrm{d}x = \int_0^T ax^2\,\mathrm{d}x
        + \int_0^T bx\,\mathrm{d}x = a\int_0^T x^2\,\mathrm{d}x
        + b\int_0^T x\,\mathrm{d}x\)

We can evaluate it as follows:

\(\left. \frac{a}{3}x^3 + \frac{b}{2}x^2 \right|_0^T\)

Notice that in this case there is no constant of integration, because we have an upper and lower bound. This is called a \textbf{definite integral}. 

If we didn't have the bounds, the \textbf{indefinite} integral would be:

\(\frac{a}{3}x^3 + \frac{b}{2}x^2+C\)

\end{document}